\documentclass[10pt,a4paper]{article}
\usepackage[margin=0.72in]{geometry}
\usepackage[T1]{fontenc}
\usepackage{lmodern,microtype,amsmath,amssymb,amsthm,booktabs}
\usepackage[colorlinks=true,urlcolor=blue,linkcolor=blue]{hyperref}
\newtheorem{theorem}{Theorem}
\newcommand{\R}{\mathbb R}
\newcommand{\F}{\mathcal F}
\newcommand{\ind}{\mathbf 1}
\newcommand{\norm}[1]{\left\lVert #1\right\rVert}
\title{A fixed interval already excludes square-root growth\\
\large Disproof of the printed asymptotic in TLMC 00000001906}
\author{lizaixi01}
\date{24 October 2026}
\begin{document}
\maketitle
\vspace{-1.3em}
\begin{abstract}
For the conventional unweighted, rough Rubio de Francia square function on the real line, every admissible family of constants satisfies
\(C_p\ge p/(8\pi e)\) for all real \(p\ge2\). A single fixed frequency interval suffices. The Lean proof constructs the native Fourier projections of genuine interval indicators, proves their integral identity by Fubini, and derives the native \(L^p\) bounds. It excludes even \(C_p=O(\sqrt p)\), hence any finite-coefficient square-root asymptotic.
\end{abstract}

\section{Precise model and scope}
The printed bilingual conjecture names the Rubio de Francia inequality for arbitrary intervals and proposes a definite-coefficient \(p^{1/2}\) asymptotic. It does not specify the underlying domain or a Fourier normalization. We use the usual real-line, unweighted, \emph{rough} interval projections, as in \cite[pp.~1--2]{grs} and \cite[p.~2, Eq.~(1.1)]{dp}:
\[
 \F f(\xi)=\int_\R f(y)e^{-2\pi i y\xi}\,dy,
 \qquad P_{[a,b]}f=\F^{-1}(\ind_{[a,b]}\F f),
 \qquad S_\Omega f=\Big(\sum_{I\in\Omega}|P_I f|^2\Big)^{1/2}.
\]
Here \(\Omega\) is any finite family of pairwise disjoint bounded intervals. Endpoints have Lebesgue measure zero. A nonnegative \(C\) is called admissible at \(p\) when the native extended \(L^p\) inequality \(\norm{S_\Omega f}_p\le C\norm f_p\) holds for all these families and all integrable \(f\in L^p(\R)\). Every usual Rubio de Francia constant necessarily satisfies this test class. The singleton family \(\{[0,1]\}\) has \(S_\Omega f=|P_{[0,1]}f|\).

We refute only the printed square-root asymptotic under this standard model. We neither dispute Rubio de Francia boundedness nor compute a sharp linear coefficient. We do not replace the operator by a smooth, martingale, or periodic square function. Theorem~1.1 of \cite{grs} concerns the lacunary operator; its later discussion separately treats arbitrary intervals on the circle. Those results motivate the scope but are not assumptions of our Lean proof. No transference or dilation theorem is needed.

\section{The actual Fourier projection}
For \(R\ge1\), put \(f_R=\ind_{[-R,R]}\). Native Lebesgue integration proves \(f_R\in L^1\cap L^p\), \(p>0\), and
\[
 \norm{f_R}_p=(2R)^{1/p}.
\]
The Fourier transform is bounded and continuous. Its restriction to \([0,1]\) is integrable, so its native inverse Fourier transform is continuous. The product phase in the iterated frequency/spatial integral has modulus one on a finite product measure. Fubini and the elementary exponential integral therefore give, whenever \(x>R\),
\[
 P_{[0,1]}f_R(x)=\int_{-R}^{R} K(x-y)\,dy,
 \qquad K(t)=\frac{e^{2\pi i t}-1}{2\pi i t}\quad(t>0).
\]
Every kernel argument here is strictly positive. Thus no totalized division at zero is used. Substitution \(x=R+\delta\), \(\delta>0\), and taking the imaginary part yield the proved identity
\begin{equation}\label{eq:imag}
 \operatorname{Im}P_{[0,1]}f_R(R+\delta)
 =\frac1{2\pi}\int_\delta^{2R+\delta}\frac{1-\cos(2\pi t)}t\,dt.
\end{equation}
The explicit integral is derived from the actual native Fourier operator; it is not substituted for its definition.

\section{A uniform oscillatory bound}
For \(1\le a\le b\), differentiate
\[
 F(t)=\frac{\sin(2\pi t)}{2\pi t},\qquad
 F'(t)=\frac{\cos(2\pi t)}t-\frac{\sin(2\pi t)}{2\pi t^2}.
\]
The fundamental theorem of calculus, \(-1\le\sin\le1\), and the actual integral of \(t^{-2}\) give
\[
 \int_a^b\frac{\cos(2\pi t)}t\,dt
 \le\frac1{2\pi b}+\frac1{2\pi a}
   +\frac1{2\pi}\Big(\frac1a-\frac1b\Big)
 =\frac1{\pi a}\le1.
\]
Only this one-sided bound is needed. If \(1\le\delta\le2\), then
\((2R+\delta)/\delta\ge R\). Combining the logarithmic integral with \eqref{eq:imag} proves
\begin{equation}\label{eq:lower}
 \operatorname{Im}P_{[0,1]}f_R(R+\delta)
 \ge\frac{\log R-1}{2\pi}.
\end{equation}

\section{Native norm bounds and the asymptotic contradiction}
Take any real \(p\ge2\) and \(R=e^p\). The genuine interval
\(E_p=[R+1,R+2]\) has Lebesgue measure one. On this entire interval, \eqref{eq:lower} and \(\operatorname{Im}z\le|z|\) imply
\[
 |P_{[0,1]}f_R(x)|\ge\frac{p-1}{2\pi},\qquad
 \norm{P_{[0,1]}f_R}_p\ge\frac{p-1}{2\pi},\qquad
 \norm{f_R}_p=2^{1/p}e\le2e.
\]
The proof of the output lower bound uses monotonicity of the native \texttt{eLpNorm} and the norm of a constant indicator on the actual measure-one interval. If a finite admissible bound is assumed, continuity and that bound explicitly imply native \texttt{MemLp} for the output. No infinite norm is converted to a default real zero.

Applying the true singleton inequality now gives
\[
 \frac{p-1}{2\pi}\le2eC_p,\qquad
 C_p\ge\frac{p-1}{4\pi e}\ge\frac{p}{8\pi e}.
\]
\begin{theorem}
For every family \(C:\R\to\R\) admissible at every real \(p\ge2\), native \(C\ne O_{p\to\infty}(\sqrt p)\). Moreover, for every finite real \(L\), the actual ratio \(C_p/\sqrt p\) does not converge to \(L\).
\end{theorem}
\begin{proof}
A supposed eventual bound \(|C_p|\le K\sqrt p\) contradicts the linear lower bound by taking \(p=t^2\) with \(t\) larger than both the threshold and \(8\pi eK\). A convergent ratio is eventually bounded, which would give precisely such a native big-O bound. Both conclusions are proved for the real \texttt{atTop} filter, so no passage from integer exponents is assumed.
\end{proof}
No optimal constant is defined by an unchecked infimum, and no norm, moment, or growth value is assigned.

\newpage
\section{Formal coverage and reproduction}
The main namespace is \texttt{TLMC1906}. The Fourier identities are proved in:
\begin{center}\begin{tabular}{@{}l@{}}
\texttt{projection\_input\_double}, \texttt{projection\_input\_kernel\_shift},\\
\texttt{projection\_input\_im}.
\end{tabular}\end{center}
The actual integral and norm estimates are proved in:
\begin{center}\begin{tabular}{@{}l@{}}
\texttt{cosine\_integral\_bound}, \texttt{exponential\_output\_lower},\\
\texttt{Admissible.output\_memLp}, \texttt{Admissible.linear\_lower}.
\end{tabular}\end{center}
The native asymptotic conclusions are:
\begin{center}\begin{tabular}{@{}l@{}}
\texttt{no\_sqrt\_bigO}, \texttt{no\_finite\_sqrt\_coefficient}.
\end{tabular}\end{center}

The project pins Lean~4.33.0 and Mathlib commit
\texttt{db584cd6d46c92f209a44c0f1c829460d327499d}.
From \texttt{lean/}, run:
\begin{verbatim}
lean --version
lake build +Main
lake env lean -DwarningAsError=true Check.lean
lake env leanchecker --fresh Main
\end{verbatim}
Equivalently, run \texttt{lean --version} followed by \texttt{./verify.sh} for the remaining three commands. The exhaustive check traverses all owned declarations, including generated ones, and invokes ordinary \texttt{\#print axioms}. Only \texttt{propext}, \texttt{Classical.choice}, and \texttt{Quot.sound} are permitted. Source-bound command, cleanup, hash and CPU-affinity receipts accompany the submission. Production timings are diagnostic observations, not controlled efficiency evidence. The complete original bilingual statement is retained in \texttt{SOURCE.md}.

\begin{thebibliography}{9}
\bibitem{grs} R.~Garg, L.~Roncal and S.~Shrivastava,
\emph{Quantitative weighted estimates for Rubio de Francia's Littlewood--Paley square function},
arXiv:1809.02937v2 (7 October 2019), pp.~1--2.
\url{https://arxiv.org/abs/1809.02937}.
\bibitem{dp} F.~Di Plinio, M.~Fl\'orez-Amatriain, I.~Parissis and L.~Roncal,
\emph{Pointwise localization and sharp weighted bounds for Rubio de Francia square functions},
arXiv:2308.01442v3 (2024), pp.~1--2, especially Eq.~(1.1).
\url{https://arxiv.org/abs/2308.01442}.
\end{thebibliography}
\end{document}
